%!TEX root=../../assign.tex

\subsection{P6}

\subsubsection{\rom{1}}
\begin{proof}
    let $\epsilon > 0$ be arbitrary in $\R$, we want to find $n_0 $ in $\N$ such that for all $n \geq n_0$ in $\N$, 
    \begin{equation*}
        \abs{x_n - 4} < \epsilon. 
    \end{equation*} 

    By Archidemean Property, there exists a $n$ in $\N$ such that 
    \begin{equation*}
        \frac{1}{n} < \frac{\epsilon}{60}. 
    \end{equation*}

    This equals to 
    \begin{equation*}
        \frac{60}{n} < \epsilon. 
    \end{equation*}

    Indeed, when $n \geq 2$, $n^3 - 6 \geq \frac{1}{4}n^3$. 

    We have 
    \begin{align*}
        \abs{\frac{4n^3 + 3n}{n^3 - 6} - 4} &= \abs{\frac{
            3n + 24
        }{
            n^3 - 6
        }} <= 
        \abs{
            \frac{
                3n + 12n
            }{n^3 - 6}
        } 
        =
        \frac{
            3n + 12n
        }{n^3 - 6} \\
        & \leq \frac{
            15n
        }{
            \frac{1}{4}n^3
        } = 
        \frac{60n}{
            n^3
        } = \frac{60}{n^2} \leq 
        \frac{
            60
        }{n
        } < \epsilon. 
    \end{align*}

    when $n \geq \max \paren{n_0, 2}$. 

    Therefore, we prove 
    \begin{equation*}
        \paren{\forall \epsilon \in \R }\paren{\exists n \in \N } \bracks{\paren{\forall n \geq n_0} \implies \paren{\abs{\frac{
            4n^3 + 3n
        }{
            n^3 - 6
        } - 4} < \epsilon}}. 
    \end{equation*}
\end{proof}

\subsubsection{\rom{2}}
\begin{proof}
    Suppose $\set{a_n}$ be a convergent sequence of integers. 

    If $a_n \neq a_{n_0}$,  
    \begin{equation*}
        \abs{a_n - a_{n_0}} \geq 1. 
    \end{equation*}

    By hypothesis, since $\set{a_n}$ is a convergent sequence, 
    \begin{equation*}
        \paren{\forall \epsilon > 0} \paren{\exists n_0 \in \N }\bracks{\paren{\forall n \geq n_0} \implies \paren{\abs{a_n - L} < \epsilon}}
    \end{equation*}

    We pick $\epsilon = \frac{1}{2}$. By definition, there exists $N$ in $\N$ such that 
    \begin{equation*}
        \paren{\forall n \geq N}\abs{x_n - L} < \frac{1}{2}. 
    \end{equation*}

    Suppose $n, m \geq N$, we have 
    \begin{equation*}
        \abs{a_n - a_m } \leq \abs{a_n - L} + \abs{a_m - L} < \frac{1}{2} + \frac{1}{2} = 1. 
    \end{equation*} 

    Since when $a_n \neq a_m$, there difference must be greater or equal to $1$. 

    Hence, $\abs{a_n - a_m} = 0$. 

    Thus, there exists a $n_0$ in $\N$ such that for all $n \geq N, a_n = a_{n_0}$. 
\end{proof}